%----------
%	VARIABLES DEL ANÁLISIS — sexismo lingüístico
%----------

El libro de códigos de IRIS define quince variables de sexismo lingüístico,
codificadas como V25 a V39, que se recogen en la
Tabla~\ref{tab:iris-inventario-variables}. Cada una captura un mecanismo
distinto mediante el cual el lenguaje periodístico puede reproducir
desigualdad.

\begin{table}[H]
    \centering
    \begingroup
    \small
    \ttabbox[\FBwidth]{
        \caption{Inventario de variables de sexismo lingüístico del libro de códigos}
        \label{tab:iris-inventario-variables}
    }{
        \begin{tabularx}{\textwidth}{
            >{\centering\arraybackslash}p{1.2cm}
            X
            >{\columncolor{gray!15}\centering\arraybackslash}p{2.2cm}
        }
            \hline
            \textbf{Cód.} & \textbf{Variable} & \textbf{En este trabajo} \\
            \hline
            V25 & \texttt{lenguaje\_sexista} & \textbf{Sí} \\
            V26 & \texttt{masc\_generico} & \textbf{Sí} \\
            V27 & \texttt{hombre\_denominar\_humanidad} & No \\
            V28 & \texttt{uso\_dual\_zorr} & No \\
            V29 & \texttt{uso\_cargo\_mujer} & No \\
            V30 & \texttt{sexismo\_discurso} & \textbf{Sí} \\
            V31 & \texttt{androcentrismo} & No \\
            V32 & \texttt{mencion\_nombre\_investigadora} & No \\
            V33 & \texttt{asimetria\_mujer\_hombre} & \textbf{Sí} \\
            V34 & \texttt{disminutivos\_infantilizacion} & No \\
            V35 & \texttt{denominacion\_sexualizada} & \textbf{Sí} \\
            V36 & \texttt{denominacion\_redundante} & No \\
            V37 & \texttt{denominacion\_dependiente} & No \\
            V38 & \texttt{criterios\_excepcion} & No \\
            V39 & \texttt{comparacion\_mujer\_hombre} & No \\
            \hline
        \end{tabularx}
    }
    \endgroup
\end{table}

\subsubsection{Selección de las cinco variables}

El alcance de este trabajo se restringe a cinco del bloque de Lenguaje (V25, V26, V30, V33 y
V35), seleccionadas por el equipo experto del proyecto atendiendo a su relevancia teórica dentro del marco de análisis. Las diez variables restantes permanecen como
marco teórico del libro de códigos, pero no se someten a evaluación automática.

Esta restricción responde asimismo a una consideración de diseño experimental:
evaluar cinco variables sobre la totalidad del subconjunto anotado y con varios
proveedores de modelos resulta preferible, en términos de solidez de las
conclusiones, a una cobertura amplia pero superficial de las quince.

\subsubsection{Definición operativa de las variables evaluadas}
\label{sec:iris_variables}

Las cinco variables seleccionadas, definidas a partir de las guías de lenguaje no
sexista y del conocimiento experto del proyecto, son las siguientes:

\begin{itemize}
    \item \textbf{V25 --- \texttt{lenguaje\_sexista}.} Uso discriminatorio del
    lenguaje por razón de sexo. Opera como categoría agregadora: engloba todo
    sexismo lingüístico de forma, de modo que la detección de cualquiera de sus
    subtipos debería implicar su marcado. Admite un tercer valor específico para
    el \emph{salto semántico}, esto es, el caso en que un masculino presentado
    como genérico se revela posteriormente no inclusivo al aparecer una marca
    explícita de mujer.

    \item \textbf{V26 --- \texttt{masc\_generico}.} Empleo del masculino
    gramatical como pretendido genérico cuando existe alternativa inclusiva.
    Constituye un subtipo específico de V25.

    \item \textbf{V30 --- \texttt{sexismo\_discurso}.} Sexismo situado en el
    \emph{fondo} del mensaje —en aquello que se cuenta— y no en la forma
    lingüística elegida. Es la única de las cinco que no constituye un subtipo
    de V25, lo que permite que una pieza presente una variable sin la otra.

    \item \textbf{V33 --- \texttt{asimetria\_mujer\_hombre}.} Tratamiento
    asimétrico de mujeres y hombres dentro de una misma pieza: diferencias en el
    uso del nombre, el apellido, el título profesional o el tratamiento.

    \item \textbf{V35 --- \texttt{denominacion\_sexualizada}.} Identificación de
    la mujer por su aspecto físico, su rol familiar o su relación con un varón,
    en detrimento de su función profesional o del motivo informativo por el que
    aparece en la pieza.
\end{itemize}

Todas las variables se codifican en una escala en la que el valor~1 indica
ausencia del fenómeno y el valor~2 su presencia; V25 admite adicionalmente el
valor~3 para el caso del salto semántico. En el tratamiento de los resultados,
y con objeto de homogeneizar la comparación entre variables, los valores 2 y 3
se agrupan como clase positiva.

\subsubsection{Prevalencia y sus implicaciones}

La distribución de las clases en la anotación experta resulta marcadamente
desigual y condiciona de forma decisiva la elección de métricas de evaluación,
tal como se analiza en la Sección~\ref{subsec:iris_divergencia}. La
Tabla~\ref{tab:iris-prevalencia-imp} recoge la prevalencia de la clase positiva de
cada variable y, junto a ella, la exactitud que alcanzaría un clasificador trivial
que respondiese siempre la clase mayoritaria.

\begin{table}[H]
    \centering
    \begingroup
    \small
    \ttabbox[\FBwidth]{
        \caption{Prevalencia de la clase positiva y exactitud del clasificador trivial ($N=1.313$)}
        \label{tab:iris-prevalencia-imp}
    }{
        \begin{tabularx}{\textwidth}{
            l
            X
            >{\centering\arraybackslash}X
            >{\columncolor{gray!15}\centering\arraybackslash}X
        }
            \hline
            \textbf{Cód.} & \textbf{Variable} & \textbf{Prevalencia «Sí»} & \textbf{Exactitud trivial} \\
            \hline
            V25 & \texttt{lenguaje\_sexista}         & 0,810 & 0,810 \\
            V26 & \texttt{masc\_generico}            & 0,810 & 0,810 \\
            V30 & \texttt{sexismo\_discurso}         & 0,428 & 0,572 \\
            V33 & \texttt{asimetria\_mujer\_hombre}  & 0,062 & 0,938 \\
            V35 & \texttt{denominacion\_sexualizada} & 0,100 & 0,900 \\
            \hline
        \end{tabularx}
    }
    \endgroup
\end{table}

La lectura de la última columna es reveladora: en V33 y V35, un modelo que no
detecte nada y responda siempre «No» acertaría el 93,8\,\% y el 90,0\,\% de las
veces. Por consiguiente, \textbf{la exactitud no constituye por sí sola una
métrica informativa} en esta tarea. Todos los resultados de este trabajo se
acompañan del coeficiente $\kappa$ de Cohen~\cite{cohen1960}, que descuenta el acuerdo esperable por
azar, de métricas F1 y de la prevalencia de cada variable.

\subsubsection{De la definición experta a la especificación ejecutable}

Para su empleo por sistemas automáticos, cada variable, partiendo de las guías de
lenguaje no sexista y del criterio experto, se formaliza en una especificación
estructurada que recoge, además de la definición conceptual: un criterio operativo
de decisión, una metodología paso a paso, ejemplos positivos y contraejemplos
comentados, casos límite documentados y las fronteras que la separan de las
variables contiguas. Esta especificación constituye la fuente
única a partir de la cual se generan tanto los \textit{prompts} del sistema de
referencia como las \textit{skills} de la arquitectura de agentes, garantizando
que ambos niveles experimentales operan con \emph{idéntico} contenido
metodológico y difieren únicamente en la forma en que dicho contenido se pone a
disposición del modelo. Los detalles de esta formalización se abordan en el
Capítulo~\ref{implementation}.
